\documentclass{ximera}
\title{Runtime Test 13: Answer Formats and Tolerance}
\begin{document}
\begin{abstract}
Tests integer, float, string, expression, and tolerance handling as isolated
answer-format behaviors.
\end{abstract}
\maketitle

\section{Integer Format}

\begin{problem}
Enter the integer \(17\).
\[
\answer[id=formatInteger,format=integer]{17}
\]
A decimal such as \(17.5\) should not be accepted as the integer \(17\).
\end{problem}

\section{Float Format}

\begin{problem}
Enter \(2.5\).
\[
\answer[id=formatFloat,format=float]{2.5}
\]
\end{problem}

\section{String Format}

\begin{problem}
Enter the word \(\text{alpha}\).
\[
\answer[id=formatString,format=string]{alpha}
\]
String comparison should follow the runtime's documented case-handling rules.
\end{problem}

\section{Expression Format}

\begin{problem}
Enter an expression equivalent to \(2(x+1)\).
\[
\answer[id=formatExpression]{2(x+1)}
\]
An algebraically equivalent response such as \(2x+2\) should be accepted.
\end{problem}

\section{Tolerance}

\begin{problem}
Enter a decimal approximation to \(\frac{1}{3}\) within tolerance \(0.01\).
\[
\answer[id=formatTolerance,format=float,tolerance=0.01]{0.3333333333}
\]
\end{problem}

Reload after correct and incorrect attempts and confirm response/correctness state
restores consistently.
\end{document}
